\documentclass[12pt]{iopart}
%\documentclass[aps,rmp,superscriptaddress,showpacs,twocolumn,longbibliography]{revtex4-1}
%\documentclass[entropy,review,submit,moreauthors,pdftex,10pt,a4paper]{mdpi} 


%Packages
\usepackage{iopams}  
\expandafter\let\csname equation*\endcsname\relax
\expandafter\let\csname endequation*\endcsname\relax

%% For landscape geometry
%\usepackage[landscape]{geometry}

% For landscape geometry
\usepackage[a4paper,margin=0.75in,landscape]{geometry}
\usepackage[fontsize=16pt]{scrextend}

\usepackage{bm}				% bold math symbols
%%\usepackage{bbm}		    % double stroked symbols
\usepackage{amsmath}
\usepackage{amssymb}
\usepackage{latexsym}
\usepackage{amsfonts}
\usepackage{graphicx}
\usepackage[usenames,dvipsnames]{xcolor}
\usepackage{dsfont}

\usepackage{color}
\definecolor{myblue}{RGB}{65,105,225}
\definecolor{mygreen}{RGB}{34,139,34}
\definecolor{myorange}{RGB}{255,69,0}
\usepackage[colorlinks,linkcolor=myorange,urlcolor=myblue,citecolor=myblue]{hyperref}

% The following packages are used to write FORTRAN code in the text
\usepackage[T1]{fontenc}
\usepackage{xcolor}
\usepackage{lmodern}
\usepackage{listings}
\lstset{language=[90]Fortran,
  basicstyle=\ttfamily,
  keywordstyle=\color{orange},
  commentstyle=\color{gray},
  morecomment=[l]{!\ }% Comment only with space after !
}

% Package to add frames: use \begin{mdframed} ... \end{mdframed}
\usepackage{mdframed}

\usepackage[spanish,activeacute]{babel}
\usepackage[]{floatflt}
\usepackage{epsfig}
\usepackage{wrapfig}
%\usepackage{html}

%\def\az{\textcolor{myblue}}
%\def\ro{\textcolor{myorange}}
\def\az{\textcolor{blue}}
\def\ro{\textcolor{red}}

\newcommand{\cA}{{\cal A}}
\newcommand{\w}{\omega}
\newcommand{\W}{\Omega}
\newcommand{\hx}{\hat{\bf x}}
\newcommand{\R}{\mathbb{R}}
\newcommand{\fx}{f_{\hx}}
\newcommand{\Fx}{F_{\hx}}
\newcommand{\hG}{\hat{\bf G}}
\newcommand{\hH}{\hat{\bf H}}
\newcommand{\hy}{\hat{\bf y}}
\newcommand{\hB}{\hat{\bf B}}
\newcommand{\hNB}{\hat{\bf N}_{B}}
\newcommand{\hNG}{\hat{\bf N}_{G}}
\newcommand{\hU}{\hat{\bf U}}
\newcommand{\hu}{\hat{\bf u}}
\newcommand{\hP}{\hat{\bf P}}
\newcommand{\hchi}{\hat{\mathbf{\chi}}}
\newcommand{\htt}{\hat{\bf t}}
\newcommand{\fy}{f_{\hy}}
\newcommand{\Fy}{F_{\hy}}
\newcommand{\hZ}{\hat{\bf Z}_n}
\newcommand{\fZ}{f_{\hZ}}
\newcommand{\hS}{\hat{\bf S}_n}
\newcommand{\hV}{\hat{\bm{\sigma}}_n}
\newcommand{\fxy}{f_{\hx\hy}}
\newcommand{\Fxy}{F_{\hx\hy}}
\newcommand{\hp}{\hat{\bf p}}
\newcommand{\hI}{\hat{\bf I}_1}
\newcommand{\hmM}{\hat{\bm \mu}_M}
\newcommand{\hsM}{\hat{\bm \sigma}_M}
\newcommand{\fxB}{f_{\hx,\hB}}

% New commands
%% Quantum mechanics notation
\newcommand{\ket}[1]{\vert{#1}\rangle} 
\newcommand{\bra}[1]{\langle{#1}\vert} 
\newcommand{\bracket}[2]{\langle{#1}\vert{#2}\rangle} 
% For operator Hilbert space
\newcommand{\kket}[1]{\vert{#1}\rangle\rangle} 
\newcommand{\bbra}[1]{\langle\langle{#1}\vert} 
\newcommand{\bbraket}[2]{\langle\langle{#1}\vert{#2}\rangle\rangle} 
% projector onto a state
\newcommand{\proj}[1]{\ket{#1}\!\bra{#1}}
%operator
\newcommand{\op}[2]{\ket{#1}\!\bra{#2}}
% expectation value
\newcommand{\mean}[1]{\langle #1 \rangle}
%identity
\newcommand{\one}{\openone}
% Trace operator
%\newcommand{\Tr}{Tr}
\newcommand{\cov}{cov}
%\DeclareMathOperator{\Tr}{Tr}
%\DeclareMathOperator{\cov}{cov}
%Fock-Liouville space
\newcommand{\opfl}[1] {\vert \vert #1 \rangle}

% general mathematical notation
\newcommand{\abs}[1]{\left|#1\right|} 
\newcommand{\pare}[1]{\left( #1 \right)}
\newcommand{\be}{\begin{equation}}
\newcommand{\ee}{\end{equation}}
\newcommand{\key}[1]{\left\{ #1 \right\}}
\newcommand{\cor}[1]{\left[ #1 \right]}
\newcommand{\acor}[1]{\left\{ #1 \right\}}
\newcommand{\la}{\langle}
\newcommand{\ra}{\rangle}
\newcommand{\bc}{\begin{center}}
\newcommand{\ec}{\end{center}}

%% Shortcuts
\newcommand{\ben}{\begin{eqnarray}}
\newcommand{\een}{\end{eqnarray}}
\newcommand{\up}{\uparrow}
\newcommand{\down}{\downarrow}

% Unitary operator U eigenvectors
\newcommand{\kpa}{\ket{\psi_\alpha}}
\newcommand{\kpb}{\ket{\psi_\beta}}

\renewcommand{\vec}[1]{\boldsymbol{#1}}
\renewcommand*{\Re}{\operatorname{Re}}

%Corrections
\newcommand{\blue}[1]{{\color{blue} #1}}

%Caligraphic Letters
\newcommand{\cHt}{{\cal H}_T}
\newcommand{\cH}{{\cal H}}
\newcommand{\cHe}{{\cal H}_E}
\newcommand{\cL}{{\cal L}}
\newcommand{\cS}{{\cal S}}
\newcommand{\cU}{{\cal U}}
\newcommand{\cW}{{\cal W}}

% Operators in the interaction picture: tilde notation
\newcommand{\tO}{\tilde{O}}
\newcommand{\trhoT}{\tilde{\rho}_T}
\newcommand{\trho}{\tilde{\rho}}
\newcommand{\tV}{\tilde{V}}
\newcommand{\tE}{\tilde{E}}

%Adjoint spaces
\newcommand{\cB}{{\cal B}}
\newcommand{\cBH}{{\cal B}({\cal H})}
\newcommand{\cBHt}{{\cal B}({\cal H}_T)}
\newcommand{\cBHe}{{\cal B}({\cal H}_E)}
\newcommand{\cBBHt}{{\cal B}({\cal B}({\cal H}_T))}
\newcommand{\cBBH}{{\cal B}({\cal B}({\cal H}))}
\newcommand{\cBa}{\cB_{\alpha\beta}}
\newcommand{\cBaa}{\cB_{\alpha\alpha}}
\newcommand{\rhoa}{\rho_{\alpha\beta}}
\newcommand{\rhoaa}{\rho_{\alpha}}
\newcommand{\id}{\mathds{1}}
\newcommand{\ide}{\mathds{1}_E}

%Text
\newcommand{\PP}{\textrm{P}}
\newcommand{\ii}{\textrm{i}}
\newcommand{\te}[1]{\textrm{#1}}

%Superoperators
\newcommand{\cLl}{{\cal L}_{\lambda}}

%Spectrum
\newcommand{\oa}{\omega_{\alpha\beta\nu}(\lambda)}
\newcommand{\loa}{\hat{\omega}_{\alpha\beta\nu}(\lambda)}
%\newcommand{\oa}{\psi_{\alpha\beta\nu}(\lambda)}
%\newcommand{\loa}{\tilde{\psi}_{\alpha\beta\nu}(\lambda)}



%%%%%%%%%%%%%%%%%%%%%%%%%%%%%%%%%%%%%%%%%%
\begin{document}

\title{\textcolor{gray}{F{\'\i}sica Computacional} \\ Resoluci{\'o}n de ecuaciones en derivadas parciales: la ecuaci{\'o}n de Schr\"odinger}


%\vspace{1.5cm}
%\hspace{0.5cm} \noindent Departamento de Electromagnetismo y F\'{\i}sica de la Materia 
%
%\hspace{0.5cm} Instituto Carlos I de F{\'\i}sica Te{\'o}rica y Computacional 
%
%\hspace{0.5cm} Universidad de Granada. E-18071 Granada. Spain

%\begin{abstract}
%These notes correspond to the first part (20 hours) of a course on Computational Methods in Nonlinear Physics within the Master on Physics and Mathematics (\emph{FisyMat}) of University of Granada. In this second chapter we describe the use of random numbers and Monte Carlo methods to perform integrals, and the adecuacy of this method for high-dimensional problems.
%\end{abstract}
%
%\vspace{2pc}
%\noindent{\it Keywords}: Computational physics, probability and statistics, Monte Carlo methods, stochastic differential equations, Langevin equation, Fokker-Planck equation, molecular dynamics.
%
%%\vspace{4cm}
%
%\section*{References and sources}
%\begin{thebibliography}{10}
%
%\bibitem{toral} R. Toral and P. Collet, \emph{Stochastic Numerical Methods}, Wiley (2014).
%
%%\bibitem{wiki} Wikipedia: https://en.wikipedia.org
%
%\end{thebibliography}


\newpage

%\tableofcontents
%
%\newpage


%\begin{figure}
%\centerline{\includegraphics[width=12cm]{}}
%\caption{}
%\label{fig1}
%\end{figure}



\section{Fundamento Te\'orico}

\begin{itemize}

\item En \ro{Mec\'anica cu\'antica}, el estado f\'{\i}sico de un sistema unidimensional de una part\'{\i}cula viene
descrito completamente por una \az{funci\'on de onda compleja $\Phi(x,t)$} que obedece la denominada \az{ecuaci\'on 
de Schr\"odinger}
%
\begin{equation}
\label{sch}
i\hbar \frac{\partial \Phi(x,t)}{\partial t} =
\left[ -\frac{\hbar^2}{2m}  \frac{\partial^2}{\partial x^2} + V(x) \right] \Phi(x,t)= H\Phi
\end{equation}
%
donde $V(x)$ es el potencial al que est\'a sometida la part\'{\i}cula, 
que supondremos que tiene una masa $m$ independiente del tiempo, $\hbar$  es la constante
de Planck reducida y \az{$H$ es el operador Hamiltoniano, que es herm\'{\i}tico}.

\item La \ro{densidad de probabilidad} de encontrar a la part\'{\i}cula en un volumen $dV$ alrededor del 
punto $x$ y en el instante $t$ es 
%
\begin{equation}
dP=|\Phi(x,t)|^2 dV.
\end{equation}
%

\item \ro{Normalizaci'on}: Para un sistema compuesto \'unicamente por una part\'{\i}cula, como es el nuestro, la
probabilidad total de encontrar la part\'{\i}cula en alg\'un punto del espacio,
en cualquier instante $t$, es igual la unidad,
%
\begin{equation}
\label{3}
\int_V |\Phi(x,t)|^2 dV=1.
\end{equation}
%

\item N\'otese que \az{la integral de probabilidad es constante durante la evoluci\'on
temporal}, lo que habr\'a que tenerse en cuenta al plantear la resoluci\'on
num\'erica del problema.

\item Podemos observar que \az{la ecuaci\'on de Schr\"odinger es lineal} y
que se trata de una \az{ecuaci\'on diferencial de primer orden en el tiempo}.

\item Para evitar el uso continuado de $m$ y $\hbar$, haremos el \ro{cambio de variable} $t\rightarrow t\hbar$
en el tiempo  y $x\rightarrow x \hbar/\sqrt{2m}$ en 
el espacio, lo que conduce a una ecuaci{\'o}n equivalente a (\ref{sch}) pero en la que $\hbar=1$ y
$m=1/2$:
%
\begin{equation}
\label{7}
-H\Phi\equiv \left[ \frac{\partial^2}{\partial x^2}- V(x) \right]\Phi(x,t) =
-i\frac{\partial \Phi}{\partial t}.
\end{equation}

\item Esta ecuaci\'on \az{puede resolverse t\'erminos de las autofunciones del hamiltoniano $H$}, 
que, recordemos, es independiente del tiempo. 

\item Llamando $u_n$ a los estados acotados y 
$u_k$ a los estados del continuo, entonces la \ro{soluci\'on dependiente del tiempo} viene
dada de forma rigurosa por 
%
\begin{equation}
\Phi(x,t)=\sum_n a_n e^{-iE_n t} u_n(x) + \int dk a(k) e^{-iE(k)t} u_k(x),
\end{equation} 
con
\begin{equation}
a_n=\langle u_n,\phi(x,0)\rangle, \qquad a(k)=\langle u_k,\phi(x,0)\rangle,
\end{equation}
%
donde \az{$\Phi(x,0)$ es el estado inicial} del sistema (y por lo tanto $a_n$ y
$a(k)$ son los coeficientes $\Phi$ en las bases $\{u_n\}$ y $\{u_k\}$, respectivamente). 

\item De seguir este procedimiento, topar\'iamos con los \ro{inconvenientes} de que:
\begin{itemize}
\item En la mayor\'{\i}a de los casos, \az{las autofunciones y los autovalores del hamiltoniano no se pueden calcular
anal\'{\i}ticamente} y han de determinarse num\'ericamente.

\item Por motivos num\'ericos, \az{el estado inicial debe restringirse necesariamente a una combinaci\'on lineal
{\em finita} de autofunciones} y, en cualquier caso, el n\'umero de \'estas debe ser el menor 
posible para reducir el trabajo computacional. Truncar la base tiene como inconveniente a\~nadido
que \az{la evoluci\'on deja de ser unitaria}.
\end{itemize}

\item Por todo ello, en lugar de seguir el procedimiento anterior \ro{integrararemos directamente la ecuaci\'on (\ref{7})}.
Pero \az{hay que ser muy cuidadosos} con el m\'etodo num\'erico porque, en general, \az{la discretizaci\'on provoca que no
se conserve la normalizaci\'on} de la funci\'on de onda (ecuaci\'on \ref{3}).

\item La \ro{soluci{\'o}n formal de la ecuaci{\'o}n de Schr\"odinger} es
%
\begin{equation}
\label{evol0}
\Phi(x,t)=e^{-i(t-t_0)H} \Phi(x,t_0)
\end{equation}
%
donde el operador \az{$e^{-itH}$ es unitario (su inverso coindice con su adjunto) por ser $H$ herm{\'\i}tico}. Esta {\'u}ltima
propiedad ser{\'a} de utilidad  a la hora de hallar la discretizaci{\'o}n m{\'a}s adecuada para resolver el problema
num\'ericamente como veremos a continuaci\'on.

\end{itemize}







\section{M{\'e}todo num{\'e}rico}

\begin{itemize}

\item \ro{Discretizaremos el espacio y el tiempo} tomando
$x_j=jh$ y $t_n=ns$, con $j=0,1,...,N$ y  $n=0,1,2,...$,
donde  \az{$h$ es el espaciado en la discretizaci\'on espacial y $s$ es el
espaciado temporal}. 

\item La \az{funci\'on de onda} viene dada en cada punto del
ret\'{\i}culo espacio-temporal por:
\begin{equation}
\Phi(x_j,t_n)\rightarrow\Phi(jh,ns)=\Phi_{j,n}\quad\quad \text{con}\quad
j=0,1,...,N \quad \text{y} \quad  n=0,1,2,... .
\end{equation}

\item Supondremos que las \ro{condiciones de contorno} para la funci{\'o}n de onda en $j=0$ y $j=N$ son las correspondientes
a la existencia de un \az{potencial infinito en esos puntos}, esto es, la densidad
de  probabilidad de encontrar la part{\'\i}cula en dichos puntos es cero.
As\'{\i}, \az{$\Phi_{0,n}=\Phi_{N,n}=0$, en cualquier instante}. 

\item Puede definirse de forma inmediata un \az{primer algoritmo} a partir de la expresi{\'o}n 
%
\begin{equation}
\Phi_{j,n+1}= e^{-i s H}\Phi_{j,n}.
\end{equation}
%
Si $s$ es muy peque{\~n}o, el operador de evoluci{\'o}n $\exp{(-isH)}$ 
puede aproximarse por su \az{desarrollo de Taylor a primer orden} en $t$. Si adem{\'a}s observamos
que el operador Hamiltoniano no es m{\'a}s que la aplicaci{\'o}n de una derivada segunda, su \az{discretizaci{\'o}n espacial} es
obvia y en total obtenemos el siguiente algoritmo
%
\begin{equation}
\Phi_{j,n+1}=\left(1-isH_D+O(sH_D)^2\right)\Phi_{j,n},
\end{equation}
%
donde 
%
\begin{equation}
\Phi_j^{\prime\prime}\equiv \frac{\partial^2 \Phi}{\partial x^2}=\frac{1}{h^2}(\Phi_{j+1}-2\Phi_j+\Phi_{j-1})+O(h^2)
\end{equation}
%
y \az{$H_D$ viene dado por}
%
\begin{equation}
H_D f_j=- \frac{1}{h^2}(f_{j+1}-2f_j+f_{j-1})+V_jf_j
\end{equation}
%
con $V_j=V(jh)$. 

\item Este algoritmo tan simple y natural adolece de un \az{grave problema}: \ro{el operador 
$(1-isH)$ {\bf no es unitario}}. Esto implica que durante la iteraci{\'o}n del algoritmo para encontrar la evoluci{\'o}n de la 
funci{\'o}n de onda,  \az{la normalizaci{\'o}n, $\sum_j h \vert\Phi_{j,n}\vert^2$, ir{\'a} variando con el tiempo}, lo que es incompatible
con la restricci\'on de normalizaci\'on a la unidad expresada en (\ref{3}) y \az{viola la interpretaci\'on de Born}.

\item As\'{\i} pues, \az{el objetivo ser{\'a} encontrar un operador evoluci{\'o}n similar al anterior pero que sea exactamente unitario}. Esto
se  consigue haciendo uso de la \ro{aproximaci\'on de Cayley} para el operador evoluci\'on temporal
%
\begin{equation}
e^{-isH}\approx \frac{1-isH_D/2}{1+isH_D/2},
\end{equation}
%
que conduce  al  \ro{algoritmo}
%
\begin{equation}
\Phi_{j,n+1}=\frac{1-isH_D/2}{1+isH_D/2}\Phi_{j,n}.
\end{equation}
%

\item N\'otese que \az{adem{\'a}s de ser unitario, este operador es exacto hasta orden $(sH_D)^2$}.

\item Para completar el algoritmo, s\'olo resta dise{\~n}ar la \az{estrategia para  utilizar eficazmente el operador $1/(1+isH_D)$}.
Para ello \az{reescribimos la anterior ecuaci{\'o}n} como
%
\begin{equation}
\label{evol}
\Phi_{j,n+1}=\left[ \frac{2}{1+isH_D/2} -1 \right]\Phi_{j,n}=\chi_{j,n}-\Phi_{j,n},
\end{equation}
%
donde
%
\begin{equation}
\chi_{j,n}\equiv \frac{2}{1+isH_D/2}\Phi_{j,n},
\end{equation}
%
o lo que es lo mismo, dada $\Phi_{j,n}$, $j=0,..,N$, $\chi_{j,n}$ es la soluci{\'o}n de la ecuaci{\'o}n
%
\begin{equation}
\left[ 1+isH_D/2  \right]\chi_{j,n}=2\Phi_{j,n},
\end{equation}
%
que en forma expl{\'\i}cita puede escribirse como
%
\begin{equation}
\chi_{j+1,n}+\left[-2+\frac{2i}{\tilde s}-\tilde V_j \right]\chi_{j,n}+\chi_{j-1,n}=\frac{4i}{\tilde s}\Phi_{j,n}
\label{kappa}
\end{equation}
%
donde $\tilde s=s/h^2$ y $\tilde V_j=h^2V_j$.

\item De esta forma, \ro{el algoritmo de evoluci{\'o}n es claro}: \az{dada una funci{\'o}n de onda en el instante $n$ para
toda posici{\'o}n $j$ se resuelve el conjunto de ecuaciones (\ref{kappa}) y se obtiene $\chi_{j,n}$ ($j=0,..,N$)}. Con esta 
soluci{\'o}n
se recurre a la ecuaci{\'o}n (\ref{evol}) para \az{obtener las nuevas $\Phi_{j,n+1}$ ($j=0,...,N$),} iter\'andose el proceso.

\item S\'olo falta por conocer c\'omo resolver ecuaciones del tipo (\ref{kappa}), es decir, \az{c\'omo invertir 
matrices tridiagonales}.

\item En este caso, hemos de resolver un \az{conjunto de ecuaciones} del tipo
%
\begin{equation}
A_j^{-}\chi_{j-1,n}+A_j^0\chi_{j,n}+A_j^+\chi_{j+1,n}=b_{j,n} \quad j=1,\ldots ,N-1
\label{recu}
\end{equation}
%
donde $A_j^-=1$, $A_j^0=-2+2i/\tilde s-\tilde V_j$, $A_j^+=1$ y $b_{j,n}=4i\Phi_{j,n}/\tilde s$.

\item \ro{Las condiciones de contorno son $\chi_0=\chi_{N}=0$} (notemos que estas condiciones de contorno implican que en (\ref{evol})
 $\Phi_{0}$ y $\Phi_{N}$ son nulas en cualquier instante).

\item Para resolver la recurrencia (\ref{recu}) \az{suponemos que su soluci{\'o}n es del tipo}
%
\begin{equation}
\chi_{j+1,n}=\alpha_{j}\chi_{j,n}+\beta_{j,n}\quad j=0,\ldots,N-1
\label{sol1}
\end{equation}
%
donde, \az{para garantizar que se cumple $\chi_{N,n}=0$, tomaremos  $\alpha_{N-1}=\beta_{N-1,n}=0$}. 

\item \az{Sustituyendo} esta expresi{\'o}n en (\ref{recu}) obtenemos
%
\begin{equation}
\chi_{j,n}=-\frac{A_j^-}{A_j^0+A_j^+\alpha_j}\chi_{j-1}+\frac{b_{j,n}-A_j^+\beta_{j,n}}{A_j^0+A_j^+\alpha_j}.
\label{sol2}
\end{equation}
%

\item Si identificamos las ecuaciones (\ref{sol1}) y (\ref{sol2}) podemos definir las \az{recurrencias para los coeficientes $\alpha$ y
$\beta$} as\'{\i}
%
\begin{equation}
\alpha_ {j-1}=-A_j ^-\gamma_j, \qquad
\beta_{j-1,n}=\gamma_j\left(b_{j,n}-A_j^+\beta_{j,n}  \right).
\label{beta}
\end{equation}
%
donde $\gamma_j^{-1}=A_j^0+A_j^+\alpha_j$. 

\item Estas ecuaciones nos dan la \az{forma de obtener todas las $\alpha$ y $\beta$
partiendo de $j=N-1$ y obteniendo en orden decreciente $\alpha_{j}$ y
$\beta_{j,n}$ con  $j=N-2, \ldots,1,0$}. 

\item Obs\'ervese que \az{$\alpha$ no depende del tiempo} y s\'olo es necesario calcularlas una vez.

\item Una vez obtenidas las $\alpha$ y $\beta$, se usa la recurrencia (\ref{sol1}) para \az{hallar las $\chi_j$ en orden de $j$
crecientes}.

\item Conocida  $\chi_{j,n}$, y con ella $\Phi_{j,n+1}$, queda discutir qu\'e \ro{funci{\'o}n de onda inicial} se usa 
y con qu\'e potencial. La funci{\'o}n de onda inicial que vamos a usar es una \az{onda plana con una amplitud gaussiana}, esto es,
%
\begin{equation}
\Phi(x,0)=e^{ik_0x}e^{-(x-x_0)^2/2\sigma^2}.
\end{equation}
%
Notemos que con esta elecci{\'o}n, la densidad de  probabilidad de encontrar inicialmente la part{\'\i}cula en un punto $x$ es
una gaussiana centrada en $x_0$ y de anchura $\sigma$. 

\item El \az{n{\'u}mero de oscilaciones completas que la funci{\'o}n de onda} tiene
sobre la red depende de $k_0$. En particular, $k_0Nh=2\pi n_{ciclos}$. En lugar de dar como par{\'a}metro inicial $k_0$, daremos
\ro{$n_{ciclos}$}. 

\item Obviamente $n_{ciclos}=0,1,....,N$. Pero f{\'\i}sicamente no tendr{\'\i}a mucho sentido que se produjese una 
oscilaci{\'o}n completa de la funci{\'o}n de onda  entre dos puntos del ret{\'\i}culo, pues querr{\'\i}a decir que la
discretizaci\'on no es suficientemente fina. As\'{\i} pues, \ro{restringiremos el par{\'a}metro $n_{ciclos}$ a los valores 
$1,...,N/4$}. \az{De esta forma, un ciclo tendr\'a 4 puntos como m{\'\i}nimo}. 

\item La posici{\'o}n media inicial y la anchura de
la gaussiana ser\'an por ejemplo $x_0=Nh/4$ y  $\sigma=Nh/16$, aunque es recomendable \az{escribir la funci\'on de 
onda inicial de manera gen\'erica y jugar con $x_0$ y $\sigma$}.
 
\item Por {\'u}ltimo, \az{el potencial que usaremos} tendr\'a una anchura $N/5$, estar\'a entrado en $N/2$ y su altura ser\'a proporcional a 
la energ{\'\i}a de la funci{\'o}n de onda incidente: $\lambda k_0^2$ (puede usarse, por ejemplo, $\lambda=0.3$).

\item \az{En resumen}, utilizando las constantes reescaladas, la funci{\'o}n de onda en el ret{\'\i}culo se tomar{\'a}
%
\begin{equation}
\Phi_{j,0}=e^{i\tilde k_0 j} e^{-8(4j-N)^2/N^2},
\end{equation}
%
donde $\tilde k_0=k_0 h=2\pi n_{ciclos}/N$ donde $n_{ciclos}=1,...,N/4$. 

\item El \az{potencial} es
\begin{equation}
\tilde V_j =V_j h^2=
\begin{cases}
0  & \text{si  $j\not\in [2N/5,3N/5]$},  \cr
\lambda\tilde k_0^2 &  \text{si $j\in  [2N/5,3N/5]$}.
\end{cases}
\end{equation}
\vspace{0.2cm}

\item S\'olo queda por fijar el \ro{par{\'a}metro $\tilde s=s/h^2$}. Puesto que la energ{\'\i}a es proporcional a $k_0^2$
y el operador din{\'a}mico discreto tiende a ser exacto en potencias de $H s$, \az{lo {\'o}ptimo es elegir
$||H||s<1$, esto es, $k_0^2 s<1$}. As\'{\i} deducimos que $\tilde s < 1/\tilde k_0^2$. En particular \ro{tomaremos $\tilde s=1/4\tilde
k_0^2$}.

\item Resumiendo, \ro{los par{\'a}metro que se han de fijar inicialmente son $N$, $n_{ciclos}$ y $\lambda$}, pues todos
los dem{\'a}s se determinan a partir de ellos. 

\item Por lo tanto, el \ro{algoritmo para resolver la ecuaci{\'o}n de Schr\"odinger unidimensional} puede esquematizarse del siguiente modo:
%
\begin{enumerate}
\item Dar los par{\'a}metros iniciales: $N$, $n_{ciclos}$ y $\lambda$. Generar $\tilde s$, $\tilde k_0$, $\tilde V_j$,
$\Phi_{j,0}$ (incluyendo las condiciones de contorno $\Phi_{0,0}=\Phi_{N,0}=0$) y $\alpha$.
\item Calcular $\beta$ utilizando la recurrencia (\ref{beta}).
\item Calcular $\chi$ a partir de (\ref{sol1}).
\item Calcular $\Phi_{j,n+1}$ de (\ref{evol}).
\item $n=n+1$, ir a al paso 2.
\end{enumerate}

\end{itemize}


%\begin{figure}
%\hspace{5cm}\includegraphics[width=8.07cm]{goldberg1.jpg}
%\includegraphics[width=8cm]{goldberg2.jpg}
%\caption{Dispersi\'on de un paquete gaussiano por un pozo cuadrado. La energ\'ia
%  media es la mitad (izquierda) o igual (derecha) a la profundidad del pozo. Los n\'umeros
%  denotan el tiempo de cada configuraci\'on en unidades arbitrarias.}
%\end{figure}

%\begin{figure}
%\includegraphics[width=8cm]{goldberg2.jpg}
%\caption{Dispersi\'on de un paquete gaussiano por un pozo cuadrado. La energ\'ia
%  promedio es igual a la profundidad del pozo. Los n\'umeros
%  denotan el tiempo de cada configuraci\'on en unidades arbitrarias.}
%\end{figure}

%\begin{figure}
%\bc
%\hspace{1cm}\includegraphics[width=8.35cm]{goldberg3.jpg}
%\includegraphics[width=8cm]{goldberg4.jpg}
%\includegraphics[width=8.07cm]{goldberg5.jpg}
%\ec
%\caption{Dispersi\'on de un paquete gaussiano por una barrera cuadrada. La
 % energ\'ia promedio es  la mitad (izquierda), igual (centro) o el doble (derecha) de la altura de la barrera. Los n\'umeros
%  denotan el tiempo de cada configuraci\'on en unidades arbitrarias.}
%\end{figure}

%\begin{figure}
%\includegraphics[width=8cm]{goldberg4.jpg}
%\caption{Dispersi\'on de un paquete gaussiano por una barrera cuadrada. La
%  energ\'ia promedio es igual a la altura de la barrera. Los n\'umeros
%  denotan el tiempo de cada configuraci\'on en unidades arbitrarias.}
%\end{figure}
%
%\begin{figure}
%\includegraphics[width=8cm]{goldberg5.jpg}
%\caption{Dispersi\'on de un paquete gaussiano por un pozo cuadrado. La energ\'ia
%  promedio es dos veces la altura de la barrera. Los n\'umeros
%  denotan el tiempo de cada configuraci\'on en unidades arbitrarias.}
%\end{figure}

\section{Problemas}
\begin{itemize}
\item{\ro{\it Obligatorio:}} Resolver la ecuaci{\'o}n de Schr\"odinger unidimensional para un potencial cuadrado. Comprobar que se conserva la norma. 

\item{\ro{\it Voluntario 1:}} {\it Estudiar el coeficiente de transmisi{\'o}n} 

El coeficiente de  transmisi{\'o}n es la probabilidad de encontrar a la part{\'\i}cula al otro lado del obst{\'a}culo para tiempos largos. En sistemas cl{\'a}sicos este coeficiente es $1$ si la energ\'ia de la part\'icula es mayor que la energ\'ia del escal\'on y $0$ si es menor. En sistemas cu\'anticos esto cambia por la acci\'on del efecto t\'unel. Para determinar si una part\'icula es reflejada o transmitida debemos colocar detectores al final del sistema. Como ya hemos visto, la probabilidad de encontrar a la part\'icula entre los puntos $x_1$ y $x_2$ viene dada por 

\be
P(x_1,x_2)=\int_{x_1}^{x_2} \left| \Phi(x) \right|^2 dx 
\ee

Suponemos que tenemos detectores finitos, actuando a derecha e izquierda de la barrera. Si estos detectores tienen un ancho de $N/5$ la probabilidad a tiempo $n$ de detectar la part\'icula a la derecha vendr\'a dado por  $P_D(n)=\sum_{j=4N/5}^N \left|  \Phi_{j,n} \right|^2$, y la probabilidad de detectarla a la izquierda ser\'ia $P_I(n)=\sum_{j=0}^{N/5} \left|  \Phi_{j,n} \right|^2$. Despu\'es de realizar el experimento $m$ veces el coeficiente de transmisi\'on se calcula como $K=m_T/m$, donde $m_T$ es el n\'umero de veces que hemos detectado la part\'icula a la derecha del potencial. 

Si la  part\'icula es detectada no hace falta continuar con la simulaci\'on, pero si no se detecta hay que proyectarla. Esto significa que a cada paso que no haya detecci\'on hay que hacer los coeficientes $\Phi_{j}=0$ para $j \in[4N/5,N]$, si hemos aplicado el detector derecho, y $\Phi_{j}=0$ para $j \in[0,N/5]$ si ha sido el izquierdo. Para garantizar la normalizaci\'on tendremos que calcular tras cada paso el valor 

\be
k=\sum_{j=0}^N  \left| \Phi_{j}\right|^2
\ee
y multiplicar todos los elementos de la funci\'on de onda de la manera $\Phi_{j}=\frac{1}{\sqrt{k}} \Phi_{j}$.

Hay que tener en cuenta que un detector que funcione a cada paso afecta fuertemente la din\'amica del sistema (comprobadlo). As\'i habr\'a que aplicar los detectores s\'olo tras un intervalo de tiempo $n_D$. Este intervalo debe ser lo suficientemente grande para dejar el sistema evoluciona, y lo suficientemente peque\~no como para evitar reflexiones en las paredes. Los resultados dependen fuertemente de este par\'ametro.

As\'i, el algoritmo quedar\'a:

\begin{enumerate}
\item Generar la funci\'on de onda inicial.
\item Evolucionar $n_D$ pasos.
\item Calcular $P_D$.
\item Generar un n\'umero aleatorio x. Si $x<P_D$ aumentar $m_T$ e ir a 1.
\item Hacer  $\Phi_{j}=0$ para $j\in[4N/5,N]$, calcular $k=\sum_{j=0}^N  \left| \Phi_{j}\right|^2$ y hacer $\Phi_{j}=\frac{1}{\sqrt{k}} \Phi_{j}$  para todo $j$.
\item Calcular $P_I$. 
\item Generar un n\'umero aleatorio x. Si $x<P_I$ e ir a 1.
\item Hacer  $\Phi_{j}=0$ para $j\in[0,N/5]$, calcular $k=\sum_{j=0}^N  \left| \Phi_{j}\right|^2$ y hacer $\Phi_{j}=\frac{1}{\sqrt{k}} \Phi_{j}$  para todo $j$.
\item Ir a 2.
\end{enumerate}

Para obtener una buena estad\'istica habr\'a que simular el sistema al menos $10^3$ veces. 

Hay que realizar como m\'inimo: 

\begin{itemize}
\item Estudiar la dependencia en $N$ de $K$ realizando simulaciones para $N=500$, $1000$, $2000$.
\item Estudiar la dependencia en $V(x)$ de $K$, usando valores $\lambda=0, 1$; $0.3$; $0.5$; $1$; $5$; $10$.
\item Comparar con resultados te\'oricos. 
\end{itemize}

Algunas cuestiones extra que se pueden calcular son las siguientes:

\begin{itemize}
\item Estudiar la din\'amica y el coeficiente de transmisi\'on en funci\'on del par\'ametro $m_T$. Discutir en base a conceptos tales como {\bf la medida} y el  {\bf colapso de la funci\'on de onda}.
\newpage
\item Calcular los valores esperados de distintos observables (posici\'on, momento, energ\'ia cin\'etica, energ\'ia total, ...) con y sin mediciones y discutir los resultados.

Nota: El valor esperado de un observable se puede calcular como 

$$\left< O \right >=\int  \phi^*(x) \hat{O}\phi(x) dx, $$ 

donde $\hat{O}$ es el operador correspondiente al observable $\left( \hat{x}=x, \hat{p}= -i\hbar \frac{\partial}{\partial x} , \dots  \right)$.  

\end{itemize}

\end{itemize}


\section*{N{\'u}meros complejos en Fortran.}

\begin{itemize}
\item{}{\it Declaraci{\'o}n:} implicit double complex ($h$)
\item{}{\it Asignaci{\'o}n de valores constantes:} $h=(1.0d0,0.3d0)$ (=1+i0.3).
\item{}{\it Asignaci{\'o}n de variables:} $h=complex(a,b)$ (=a+ib).
\end{itemize}




%%%%%%%%%%%%%%%%%%%%%%%%%%%%%%%%%%%%%%%%%%%%%%%%%%%%%%%%%%%%%%%%%%%%%%%%%%%
\end{document}
